\section{Introduction}




Reinforcement Learning has played an integral role in enhancing the reasoning capabilities of large language models (LLMs) \citep{jaech2024openai, team2025kimi, guo2025deepseek}.
RL allows the model to freely explore the reasoning space in the post-training process and improve its performance based on the feedback provided in the environment. 
However, applying Reinforcement Learning directly to a base model (i.e., ``Zero RL'') \citep{zeng2025simplerlzooinvestigatingtamingzero} presupposes a certain level of inherent capability. This method often falters when applied to weaker models or tasks of high complexity, as the exploration process may fail to explore and discover meaningful reward signals. 
Conversely, the classical Supervised Fine-Tuning (SFT) \citep{wei2021finetuned} offers a direct and efficient method to distill knowledge from high-quality, human-annotated data, enabling models to rapidly and accurately fit the target distribution. Yet this approach often curtails the model's exploratory capabilities, potentially leading to overfitting on the demonstration data and compromising its generalization performance on out-of-distribution inputs. 
Consequently, a sequential ``SFT-then-RL'' pipeline \citep{yoshihara2025practical} has emerged as the standard, adopted by numerous state-of-the-art open-source models. While effective, this multi-stage process, which first elevates the model's capabilities through SFT before refining them with RL, is notoriously resource-intensive and usually requires careful tuning to ensure effectiveness.


To circumvent these challenges, recent works have focused on integrating SFT or SFT-style imitation learning losses directly with RL objectives \citep{yan2025learning, fu2025srft, zhang2025on}.
In these approaches, the model is updated using a composite loss function. The balance between the imitation and exploration components is governed by various strategies, including a fixed coefficient, a predefined schedule, a dynamic adjustment based on entropy, or a learnable parameter. 
These works predominantly treat the SFT and RL losses as two distinct objectives. And a detailed analysis of \textit{why these two learning signals can be effectively combined within a unified optimization process} remains largely unexplored.

Despite their distinct mathematical formulations, we find that the gradient calculations from these approaches can be viewed as a single, unified form. Inspired by Generalized Advantage Estimator \citep{schulman2015high}, we introduce \framework (UPGE), a framework that formally subsumes the gradients of various post-training objectives into one generalized expression. We provide analysis to show that the various forms of gradients are, in fact, not conflicting.
Instead, they act as complementary learning signals that can jointly guide the optimization process.
However, these gradient estimators possess different characteristics, and there exists a bias-variance tradeoff in their respective gradient components. 
Building upon this unified perspective, we propose \method (\mot), a hybrid algorithm to dynamically choose more desirable training signals by adapting a mixing ratio between the SFT and RL losses.
This mechanism allows \mot to be intrinsically adaptive to models of varying capabilities and data of differing complexities.

We implement a simple instance of \mot, which adaptively switches between SFT and RL based on rollout accuracy, and empirically demonstrate that it achieves strong results.
Our empirical evaluations demonstrate that \mot surpasses strong baselines such as SFT$\rightarrow$GRPO and LUFFY with Qwen2.5-Math-7B, achieving a 7-point gain over our strongest baseline on AIME 2024.
Moreover, \mot also yields substantial improvements even on relatively smaller and weaker models, including Qwen2.5-Math-1.5B and Llama3.1-8B.
Through detailed training dynamics and illustrative training visualizations, we clearly reveal the features and underlying mechanisms of \mot.
The following are several key takeaways:


\begin{tcolorbox}[takeawaysbox]
\begin{enumerate}[leftmargin=1em]
    \item \gpe provides a theoretical unification of a wide spectrum of post-training algorithms, covering both SFT and RL losses within a single formulation~\textbf{($\S$~\ref{sec:unified-view})}.
    \item \mot is capable of outperforming previous post-training and mixed-policy algorithms across a diverse range of models~\textbf{($\S$~\ref{sec:experiments})}.
    \item Dynamic integration of SFT and RL in \mot achieves the highest \emph{Pass@1024}, facilitating enhanced exploration and generalization of the model~\textbf{($\S$~\ref{sec:explor_exploit})}.
\end{enumerate}



\end{tcolorbox}









